%% Manuscript for Computers & Graphics (Elsevier), built on elsarticle (numbered references).
%% Journal guide: https://www.sciencedirect.com/journal/computers-and-graphics/publish/guide-for-authors
%%
%% Review is DOUBLE ANONYMIZED. This file produces the anonymized manuscript:
%%   - no author name, affiliation, ORCID, e-mail, acknowledgements or funding in the body
%%   - cite your own earlier work in the third person
%% The identifying title page lives in title_page.tex (upload it as a separate file).
%% Set \anonfalse below only to produce a non-anonymous version (e.g. arXiv preprint).
\documentclass[preprint,12pt]{elsarticle}
%% Final journal layouts (not needed for submission): [final,1p,times], [final,3p,times], [final,5p,times]

\usepackage{amssymb}
\usepackage{amsmath}
\usepackage{graphicx}
\usepackage{booktabs}
\usepackage{longtable}
%% \usepackage{lineno}   %% uncomment, plus \linenumbers after \end{frontmatter}, for line numbers

\newif\ifanon
\anontrue   %% anonymized manuscript (journal submission)
%% \anonfalse %% identified version (arXiv preprint)

\journal{Computers \& Graphics}

\begin{document}

\begin{frontmatter}

%% Title: concise, informative, no abbreviations/formulae if avoidable.
\title{}

\ifanon
\else
  \author{Paritosh Sharma}
  \ead{paritosh.sharmas@gmail.com}
  %% ORCID: https://orcid.org/0000-0001-9938-008X
  %% Independent researcher -- fill in your postal address (the guide requires a full postal address and country)
  \affiliation{organization={Independent},%
              addressline={97 chemin du fournas},
              city={Vallauris},
              postcode={06220},
              state={},
              country={France}}
\fi

%% Abstract: at most 250 words, stand-alone, no references, no uncommon abbreviations.
%% Purpose, basic procedures, main findings, principal conclusions, what is new.
\begin{abstract}
\end{abstract}

%% Graphical abstract and Highlights are submitted as SEPARATE files in the online system,
%% not in this PDF:
%%   - Graphical abstract: 531 x 1328 px (h x w) or proportional; readable at 5 x 13 cm; TIFF/EPS/PDF/MS Office.
%%   - Highlights: separate editable file with "highlights" in its name; 3 to 5 bullets, max 85 characters each (incl. spaces).

%% Keywords: 1 to 7, English, avoid multi-word keywords with "and"/"of".
\begin{keyword}
\end{keyword}

\end{frontmatter}

%% Sections are numbered 1, 1.1, 1.1.1. Cross-reference by number (\ref), never "the text".

\section{Introduction}
\label{sec:intro}

\section{Related Work}
\label{sec:related}

\subsection{Pedestrians in the loop and VR experiments}
\label{sec:rw-ped}

\subsection{Low-vision simulation and rehabilitation in VR}
\label{sec:rw-lowvision}

\subsection{Scenario specification languages and ontologies}
\label{sec:rw-spec}

\subsection{LLM-based scenario generation}
\label{sec:rw-llm}

\subsection{Pedestrian-focused scenario generation for autonomous-vehicle testing}
\label{sec:rw-pedav}

\subsection{Urban and 3D scene generation}
\label{sec:rw-urban}

\subsection{Positioning and research gap}
\label{sec:rw-gap}

\section{Problem Statement and Research Questions}
\label{sec:problem}

\section{Pedestrian Navigation Task Taxonomy and Representation}
\label{sec:taxonomy}

\subsection{Scope and inclusion criteria}
\subsection{Task categories}
\subsection{Task ontology: targets and failure conditions}
\subsection{Encoding in OpenSCENARIO and OpenDRIVE}

\section{System Overview}
\label{sec:system}

\subsection{Map ingestion and the metric map model}
\subsection{LLM agent and grounded tools}
\subsection{Scenario DSL and deterministic storyboard assembly}
\subsection{Sandboxed execution}
\subsection{Map-aware validation and repair loop}
\subsection{Interactive player (VR and first person)}

\section{Evaluation}
\label{sec:eval}

\subsection{Experimental setup}
\subsection{Benchmark: validity and prompt adherence (RQ2, RQ3)}

\begin{longtable}{@{}p{1.7cm}p{2.9cm}p{4.4cm}p{3.1cm}@{}}
\caption{Performance of the L2 pipeline per scenario type on the Sophia Antipolis map. Result: count (percentage; 95\% Wilson interval) over all runs of that type.}\label{tab:l2_results}\\
\toprule
Scenario type & Task description & Quantitative metric & Result \\
\midrule
\endfirsthead
\toprule
Scenario type & Task description & Quantitative metric & Result \\
\midrule
\endhead
\texttt{ptp\_env} (9 runs) & Reach a landmark (bus stop, bench, waste basket) over walkable ground & Accepted (validator passes the final scenario) & 9/9 (100\%; 70--100) \\
 &  & First try (accepted with no repair loop) & 3/9 (33\%; 12--65) \\
 &  & Adherent (constraint axes and operators match the request) & 9/9 (100\%; 70--100) \\
 &  & Solvable (the nominal walker satisfies all tasks) & 9/9 (100\%; 70--100) \\
 &  & Script submissions per run (mean) & 1.33 \\
 &  & Model steps per run (mean) & 3.7 \\
 &  & Wall-clock per run, s (mean) & 44 \\
\midrule
\texttt{ptp\_road} (9 runs) & Reach a spot on the carriageway & Accepted (validator passes the final scenario) & 6/9 (67\%; 35--88) \\
 &  & First try (accepted with no repair loop) & 4/9 (44\%; 19--73) \\
 &  & Adherent (constraint axes and operators match the request) & 5/9 (56\%; 27--81) \\
 &  & Solvable (the nominal walker satisfies all tasks) & 5/9 (56\%; 27--81) \\
 &  & Script submissions per run (mean) & 0.89 \\
 &  & Model steps per run (mean) & 6.4 \\
 &  & Wall-clock per run, s (mean) & 52 \\
\midrule
\texttt{cpath} (9 runs) & Reach a goal while staying inside a corridor (path or sidewalk) & Accepted (validator passes the final scenario) & 1/9 (11\%; 2--44) \\
 &  & First try (accepted with no repair loop) & 0/9 (0\%; 0--30) \\
 &  & Adherent (constraint axes and operators match the request) & 1/9 (11\%; 2--44) \\
 &  & Solvable (the nominal walker satisfies all tasks) & 1/9 (11\%; 2--44) \\
 &  & Script submissions per run (mean) & 4.89 \\
 &  & Model steps per run (mean) & 6.1 \\
 &  & Wall-clock per run, s (mean) & 121 \\
\midrule
\texttt{env\_con} (9 runs) & Reach a goal without touching static obstacles (parked cars) & Accepted (validator passes the final scenario) & 4/9 (44\%; 19--73) \\
 &  & First try (accepted with no repair loop) & 2/9 (22\%; 6--55) \\
 &  & Adherent (constraint axes and operators match the request) & 4/9 (44\%; 19--73) \\
 &  & Solvable (the nominal walker satisfies all tasks) & 4/9 (44\%; 19--73) \\
 &  & Script submissions per run (mean) & 3.11 \\
 &  & Model steps per run (mean) & 6.8 \\
 &  & Wall-clock per run, s (mean) & 185 \\
\midrule
\texttt{dyn\_nav} (9 runs) & Cross while keeping clear of a moving vehicle & Accepted (validator passes the final scenario) & 5/9 (56\%; 27--81) \\
 &  & First try (accepted with no repair loop) & 2/9 (22\%; 6--55) \\
 &  & Adherent (constraint axes and operators match the request) & 5/9 (56\%; 27--81) \\
 &  & Solvable (the nominal walker satisfies all tasks) & 4/9 (44\%; 19--73) \\
 &  & Script submissions per run (mean) & 4.33 \\
 &  & Model steps per run (mean) & 5.6 \\
 &  & Wall-clock per run, s (mean) & 185 \\
\midrule
\texttt{dep\_comp} (9 runs) & Several goals that must be reached in a given order & Accepted (validator passes the final scenario) & 5/9 (56\%; 27--81) \\
 &  & First try (accepted with no repair loop) & 0/9 (0\%; 0--30) \\
 &  & Adherent (constraint axes and operators match the request) & 5/9 (56\%; 27--81) \\
 &  & Solvable (the nominal walker satisfies all tasks) & 5/9 (56\%; 27--81) \\
 &  & Script submissions per run (mean) & 4.44 \\
 &  & Model steps per run (mean) & 8.2 \\
 &  & Wall-clock per run, s (mean) & 192 \\
\midrule
\texttt{indep\_comp} (9 runs) & Several goals that may be reached in any order & Accepted (validator passes the final scenario) & 7/9 (78\%; 45--94) \\
 &  & First try (accepted with no repair loop) & 0/9 (0\%; 0--30) \\
 &  & Adherent (constraint axes and operators match the request) & 0/9 (0\%; 0--30) \\
 &  & Solvable (the nominal walker satisfies all tasks) & 7/9 (78\%; 45--94) \\
 &  & Script submissions per run (mean) & 2.22 \\
 &  & Model steps per run (mean) & 10.2 \\
 &  & Wall-clock per run, s (mean) & 139 \\
\midrule
\texttt{mixed} (12 runs) & Combined constraints: signal compliance, a moving car, looking at it & Accepted (validator passes the final scenario) & 2/12 (17\%; 5--45) \\
 &  & First try (accepted with no repair loop) & 0/12 (0\%; 0--24) \\
 &  & Adherent (constraint axes and operators match the request) & 2/12 (17\%; 5--45) \\
 &  & Solvable (the nominal walker satisfies all tasks) & 0/12 (0\%; 0--24) \\
 &  & Script submissions per run (mean) & 2.58 \\
 &  & Model steps per run (mean) & 6.9 \\
 &  & Wall-clock per run, s (mean) & 224 \\
\bottomrule
\end{longtable}

\subsection{Vision ablation}
\subsection{Coverage analysis of the taxonomy (RQ1)}
\subsection{Qualitative feedback round (RQ4)}

\section{Discussion}
\label{sec:discussion}

\subsection{Limitations}

\section{Conclusion}
\label{sec:conclusion}

%% ------------------------------------------------------------------
%% Back matter (journal requirements). Unnumbered sections.
%% Acknowledgements are NOT placed here: for double-anonymized review they go on the title page only.
%% ------------------------------------------------------------------

\section*{CRediT authorship contribution statement}

\section*{Declaration of competing interest}

\section*{Funding}

\section*{Data availability}
%% Journal uses "Option C": deposit research data/code in a repository and cite it, or explain why not.
%% For review, link anonymously (no identifying repository URL).

%% Add this section (before the references) only if generative AI tools were used in manuscript preparation.
%% Required wording: tool name, reason, and that the author reviewed and edited the content.
\section*{Declaration of generative AI and AI-assisted technologies in the manuscript preparation process}

%% Appendices: A, B, ...; equations Eq.(A.1), tables Table A.1, figures Fig. A.1
\appendix
\section{Evaluation prompts}
\label{app:prompts}
\begin{longtable}{@{}p{3.5cm}p{2.5cm}p{6.2cm}@{}}
\caption{Evaluation prompts (Sophia Antipolis map, all requested object types exist on it).}\label{tab:prompts}\\
\toprule
Id & Type & Prompt \\
\midrule
\endhead
\texttt{sa\_ptp\_env\_1} & \texttt{ptp\_env} & Walk along the sidewalk to the nearest bus stop. \\
\texttt{sa\_ptp\_env\_2} & \texttt{ptp\_env} & The participant should find the bench nearby and walk up to it. \\
\texttt{sa\_ptp\_env\_3} & \texttt{ptp\_env} & Go to the waste basket and stop next to it. \\
\texttt{sa\_ptp\_road\_1} & \texttt{ptp\_road} & Walk out onto the road and stand on the marked spot just past a crossing. \\
\texttt{sa\_ptp\_road\_2} & \texttt{ptp\_road} & The participant must stand in the middle of the carriageway for a moment. \\
\texttt{sa\_ptp\_road\_3} & \texttt{ptp\_road} & Have the user walk onto the road surface at the far end of the street. \\
\texttt{sa\_cpath\_1} & \texttt{cpath} & Cross the road but stay inside the marked corridor along the crosswalk. \\
\texttt{sa\_cpath\_2} & \texttt{cpath} & Follow the marked path along the street without leaving it. \\
\texttt{sa\_cpath\_3} & \texttt{cpath} & Walk to the crossing using only the sidewalk; stepping off it means failure. \\
\texttt{sa\_env\_con\_1} & \texttt{env\_con} & Cross the road at the zebra crossing without touching the parked cars. \\
\texttt{sa\_env\_con\_2} & \texttt{env\_con} & Walk between the parked cars along the kerb and reach the other end without touching any of them. \\
\texttt{sa\_env\_con\_3} & \texttt{env\_con} & Avoid the parked vehicles along the street while you walk to the bus stop. \\
\texttt{sa\_dyn\_nav\_1} & \texttt{dyn\_nav} & Cross at the zebra crossing while a car approaches; avoid being hit. \\
\texttt{sa\_dyn\_nav\_2} & \texttt{dyn\_nav} & A bus drives along the street. Walk across the road without being hit by it. \\
\texttt{sa\_dyn\_nav\_3} & \texttt{dyn\_nav} & Cross the road with oncoming traffic and keep clear of the speeding car. \\
\texttt{sa\_dep\_comp\_1} & \texttt{dep\_comp} & First walk to the crossing, then cross the road, then continue to the spot on the other side. \\
\texttt{sa\_dep\_comp\_2} & \texttt{dep\_comp} & Do three things in order: go to the bus stop, then to the bench, then to the waste basket. \\
\texttt{sa\_dep\_comp\_3} & \texttt{dep\_comp} & Walk to the kerb, afterwards cross the road, and finally walk to the waste basket. \\
\texttt{sa\_indep\_comp\_1} & \texttt{indep\_comp} & Visit the waste basket and the bench, in any order. \\
\texttt{sa\_indep\_comp\_2} & \texttt{indep\_comp} & Find both the bench and the bus stop; the order does not matter. \\
\texttt{sa\_indep\_comp\_3} & \texttt{indep\_comp} & Reach a bench and a tree, either one first. \\
\texttt{sa\_mixed\_1} & \texttt{mixed} & Wait for the pedestrian light to turn green, then cross the road. \\
\texttt{sa\_mixed\_2} & \texttt{mixed} & Cross at the signalised crossing on green while a car comes down the street, and stay clear of the car. \\
\texttt{sa\_mixed\_3} & \texttt{mixed} & Look at the approaching car, then cross the road once it has passed. \\
\texttt{sa\_mixed\_4} & \texttt{mixed} & At the crossing wait until the light is green, look at the traffic light, and then cross. \\
\bottomrule
\end{longtable}


%% References: numbered, in order of first appearance; more than 6 authors -> first 6 then et al.; add DOIs.
%% Mark preprints ("preprint" / server name). Every reference must be a real source.
%% Placeholder so BibTeX/natbib do not warn while no citation exists yet. DELETE the next line once the text cites something.
\nocite{wu2022gust3d}
\bibliographystyle{elsarticle-num}
\bibliography{refs}

\end{document}
